% 第6章：目标项分解与变量分块。网格表示数据矩阵 A：n=8 行对应 8 个目标项 ell_i，d=6 列对应坐标 x_1..x_6；
% 顶部一行表示决策向量 x。左：本轮只读取第 2、6 行（批量 B={2,6}），更新全部坐标；
% 右：读取全部行，只更新第 3、4 列对应的坐标块 J={3,4}。
\begin{tikzpicture}[x=4.2mm,y=4.2mm,note/.style={font=\FigureNoteFont\fontsize{7.5}{9}\selectfont,align=center}]
  % ---------- left panel
  \begin{scope}
    \foreach \i in {2,6} {\fill[FigureBlueFill] (0,8-\i) rectangle (6,9-\i);}
    \draw[FigureGrid,line width=.4pt,step=1] (0,0) grid (6,8);
    \foreach \i in {2,6} {\draw[FigureBlue,line width=.9pt] (0,8-\i) rectangle (6,9-\i);}
    \draw[FigureStroke,line width=.6pt] (0,0) rectangle (6,8);
    \foreach \i in {2,6} {\node[left,text=FigureBlue] at (0,8.5-\i) {$\ell_{\i}$};}
    \node[left,text=FigureMuted] at (0,7.5) {$\ell_{1}$};
    \node[left,text=FigureMuted] at (0,0.5) {$\ell_{8}$};
    \fill[FigureBlueFill] (0,9) rectangle (6,10);
    \draw[FigureBlue,line width=.9pt,step=1] (0,9) grid (6,10);
    \draw[FigureBlue,line width=.9pt] (0,9) rectangle (6,10);
    \node[left] at (0,9.5) {$x$};
    \node[note] at (3,10.8) {更新全部坐标};
    \node[note] at (3,-0.9) {目标项分解：只读取 $i\in B=\{2,6\}$};
  \end{scope}
  % ---------- right panel
  \begin{scope}[xshift=13*4.2mm]
    \fill[FigureGreenFill] (2,0) rectangle (4,8);
    \draw[FigureGrid,line width=.4pt,step=1] (0,0) grid (6,8);
    \draw[FigureGreen,line width=.9pt] (2,0) rectangle (4,8);
    \draw[FigureStroke,line width=.6pt] (0,0) rectangle (6,8);
    \node[left,text=FigureMuted] at (0,7.5) {$\ell_{1}$};
    \node[left,text=FigureMuted] at (0,0.5) {$\ell_{8}$};
    \draw[FigureGrid,line width=.6pt,step=1] (0,9) grid (6,10);
    \draw[FigureStroke,line width=.6pt] (0,9) rectangle (6,10);
    \fill[FigureGreenFill] (2,9) rectangle (4,10);
    \draw[FigureGreen,line width=.9pt] (2,9) rectangle (4,10);
    \node[left] at (0,9.5) {$x$};
    \node[note] at (3,10.8) {只更新坐标块 $x_J$};
    \node[note] at (3,-0.9) {变量分块：读取全部项，$J=\{3,4\}$};
  \end{scope}
  \node[note,text=FigureMuted] at (9.5,4) {两者可以\\同时使用};
\end{tikzpicture}
